\section{Artsfredningsbekendtgørelsen}

\begin{frame}
  \vspace{3em}
  \begin{center}
    {\LARGE \textbf{Artsfredningsbekendtgørelsen}}
    \vspace{2em}
    \\ Bekendtgørelse om fredning af visse dyre- og plantearter og pleje af tilskadekommet vildt
\par
% Våbenbekendtgørelsen
{\normalsize \href{https://www.retsinformation.dk/eli/lta/2021/521}{BEK nr 521 af 25/03/2021 \\ §§ 1-9 | § 10, stk. 1 | § 11 | §§ 25-31}}
\vspace{1em}
\par
  \end{center}
\end{frame}

\begin{frame}{Anvendelsesområde}
\begin{itemize}
  \item Bekendtgørelsen omfatter:
  \begin{itemize}
    \item Fuglearter i vild tilstand i EU-Europa (fuglebeskyttelsesdirektivet)
    \item Pattedyr, krybdyr, fisk, padder, invertebrater (bilag 1)
    \item Planter (bilag 2)
    \item Dyr fra bilag 3 (igler/snegle mv.)
  \end{itemize}
  \item Gælder også udsatte/undslupne pattedyr \& fugle med vildtlevende, reproducerende bestande (ikke traditionelle husdyr)
  \item Omfatter alle livsstadier + dele/produkter (dyr \& planter)
  \item Vildt = naturligt forekommende pattedyr \& fugle i DK (inkl. trækfugle)
\end{itemize}
\end{frame}

\begin{frame}{Undtagelser}
\begin{itemize}
  \item \textbf{Gælder ikke:}
  \begin{itemize}
    \item Landbrugsmæssigt opdrættede dyr (jf. lov om hold af dyr)
    \item Dyr opdrættet i fangenskab af lovligt erhvervede forældre
  \end{itemize}
  \end{itemize}
\end{frame}

\begin{frame}{Forbud mod drab \& indfangning}
\begin{itemize}
  \item Forsætligt drab eller indfangning af vilde fugle (§ 1, stk. 1) forbudt – uanset metode
  \item Også forbud mod besiddelse af levende vilde fugle
  \item Gælder IKKE for arter der må
  \begin{itemize}
    \item jages/indfanges efter jagttidsbekendtgørelsen
    \item reguleres efter vildtskadebekendtgørelsen
  \end{itemize}
\end{itemize}
\end{frame}

\begin{frame}{Forbud mod handel mv.}
\begin{itemize}
  \item Forbud mod handel, transport til handel, opbevaring til handel, udbud til salg, konservering, udstilling af vilde fugle (§ 1, stk. 1 \& 3)
  \item Gælder også genkendelige dele \& produkter
  \item \textbf{Undtagelse:} Gælder IKKE for fugle på \textbf{bilag 4}, hvis lovligt dræbt/indfanget/erhvervet
\end{itemize}
\end{frame}

\begin{frame}{Bilag 4 – Arter der kan handles (undtagelse fra forbud)}
\tiny
Fugle fra fuglebeskyttelsesdirektivets bilag III A/B – handel tilladt hvis lovligt erhvervet.\\
* = Invasive (EU 1143/2014: forbud mod handel med levende eksemplarer)

\begin{itemize}
  \item Agerhøne \textit{(Perdix perdix)}
  \item Amerikansk skarveand \textit{(Oxyura jamaicensis)} *
  \item Canadagås \textit{(Branta canadensis)}
  \item Fasan \textit{(Phasianus colchicus)}
  \item Grågås \textit{(Anser anser)}
  \item Gråand \textit{(Anas platyrhynchos)}
  \item Krikand \textit{(Anas crecca)}
  \item Nilgås \textit{(Alopochen aegyptiaca)} *
  \item Ringdue \textit{(Columba palumbus)}
  \item Skovsneppe \textit{(Scolopax rusticola)}
\end{itemize}

\tiny
Undtagelse: Tillader handel, transport, opbevaring, konservering, udstilling mv.
\end{frame}

\begin{frame}{Beskyttelse af redetræer \& reder}
\begin{itemize}
  \item Kolonirugende fugles redetræer: fædningsforbud 1. feb.–31. jul.
  \item Rovfugle \& uglers redetræer: fældningsforbud 1. feb.–31. aug.
  \item Ørne, sort stork, rød glente: redetræer må \textbf{aldrig} fældes
  \item Hule træer \& træer med spættehuller: fældningsforbud 1. nov.–31. aug.
  \item Digesvalereder: må ikke ødelægges 1. apr.–31. aug.
\end{itemize}
\end{frame}

\begin{frame}{Import \& handel med levende fugle}
\begin{itemize}
  \item Forbud mod import af levende fugle (fuglebeskyttelsesdirektivet, artikel 1) fra lande uden for EU – medmindre dokumenteret fangenskabsopdræt
  \item Gælder også underarter (selvom ikke naturlige i EU)
  \item Handel/forsendelse af levende fugle: kun tilladt hvis fangenskabsopdrættet er registreret
  \item Dokumentationspligt over for Miljøstyrelsen
\end{itemize}
\end{frame}

\begin{frame}{Streng beskyttelse (Habitatdirektiv bilag IV)}
\begin{itemize}
  \item Gælder vildtlevende dyr på habitatdirektivets bilag IV (jf. bilag 1 i bekendtgørelsen)
  \item \textbf{Totalt forbud} mod:
  \begin{itemize}
    \item Alle former for forsætlig indfangning eller drab
    \item Forsætlig ødelæggelse/ind samling af æg i naturen
  \end{itemize}
  \item Planter på bilag IV (jf. bilag 2): Forbud mod forsætlig plukning, indsamling, afskæring, opgravning, oprivning med rod eller ødelæggelse
  \item Eksempler: Mange flagermus, odder, ulv (tidl.), visse padder/insekter, orkideer mv.
  \item \textbf{Hovedregel:} Ingen undtagelse for jagt – strengeste niveau!
\end{itemize}
\end{frame}


\begin{frame}{Handel med jagtbart vildt}
\begin{itemize}
  \item
  \begin{itemize}
    \item Handel/transport af \textbf{døde pattedyr} tilladt – hvis lovligt nedlagt/ombragt
    \item \textbf{Levende vildt:} Kun efter tilladelse fra Miljøstyrelsen
  \end{itemize}
  \item Erhvervsmæssige købere skal dokumentere køb de sidste 2 år (hvad \& fra hvem) på Miljøstyrelsens forlangende
  \item \textbf{Budskab:} Sporbarhed er vigtig – ulovligt vildt må ikke handles
\end{itemize}
\end{frame}

\begin{frame}{Bilag 5 – Forbud mod konservering (dansk natur)}
\tiny
Pattedyr og fugle – ingen udstopning/konservering uden dispensation

\begin{columns}[T]
\column{0.33\textwidth}
\textbf{Pattedyr}
\begin{itemize}
  \item Birkemus (\textit{Sicista betulina})
  \item Hasselmus (\textit{Muscardinus avellanarius})
  \item Hvaler (Cetacea, alle)
  \item Odder (\textit{Lutra lutra})
  \item Småflagermus (Microchiroptera spp.)
  \item Ulv (\textit{Canis lupus})
  \item Bæver (\textit{Castor fiber})
\end{itemize}

\column{0.33\textwidth}
\textbf{Rovfugle \& ugler}
\begin{itemize}
  \item Blå kærhøg (\textit{Circus cyaneus})
  \item Dværgfalk (\textit{Falco columbarius})
  \item Fiskeørn (\textit{Pandion haliaetus})
  \item Havørn (\textit{Haliaeetus albicilla})
  \item Hedehøg (\textit{Circus pygargus})
  \item Jagtfalk (\textit{Falco rusticolus})
  \item Kongeørn (\textit{Aquila chrysaetos})
  \item Kirkeugle (\textit{Athene noctua})
  \item Mosehornugle (\textit{Asio flammeus})
  \item Lærkefalk (\textit{Falco subbuteo})
  \item Rød glente (\textit{Milvus milvus})
  \item Stor hornugle (\textit{Bubo bubo})
  \item Slangeørn (\textit{Circaetus gallicus})
  \item Vandrefalk (\textit{Falco peregrinus})
\end{itemize}

\column{0.33\textwidth}
\textbf{Andre fugle}
\begin{itemize}
  \item Biæder (\textit{Merops apiaster})
  \item Ellekrage (\textit{Coracias garrulus})
  \item Dværgterne (\textit{Sterna albifrons})
  \item Hvid stork (\textit{Ciconia ciconia})
  \item Hærfugl (\textit{Upupa epops})
  \item Høgesanger (\textit{Sylvia nisoria})
  \item Rørdrum (\textit{Botaurus stellaris})
  \item Sandterne (\textit{Gelochelidon nilotica})
  \item Sortterne (\textit{Chlidonias niger})
  \item Skestork (\textit{Platalea leucorodia})
  \item Sort stork (\textit{Ciconia nigra})
  \item Tinksmed (\textit{Tringa glareola})
  \item Trane (\textit{Grus grus})
  \item Pirol (\textit{Oriolus oriolus})
\end{itemize}
\end{columns}
\tiny
Gælder kun dyr fra dansk natur. Dispensation mulig (typisk kun museer mv.).
\end{frame}

\begin{frame}{Pleje af almindeligt tilskadekommet vildt udført af privat}
\begin{itemize}
  \item Gælder \textbf{ikke} vildtarter anført i bilag 6
  \item Må plejes med henblik på genudsætning (uanset fredningsregler i kap. 2-4)
  \item Genudsæt hurtigst muligt – helst på oprindeligt levested
  \item Hvis ingen realistisk genudsætning inden \textbf{6 måneder} → afliv hurtigst muligt
  \item Miljøstyrelsen kan dispensere i særlige tilfælde (på vilkår)
  \item \textbf{Budskab:} Hjælp OK, men begrænset tid – ingen permanent hold
\end{itemize}
\end{frame}

\begin{frame}{Pleje af strengt beskyttede arter (bilag 6) - kun særligt bemyndigede}
\begin{itemize}
  \item 
  \begin{itemize}
    \item Bilag 6-arter (fx odder, visse flagermus, ulv, bæver, sjældne fugle): Kun pleje hos \textbf{bemyndigede} personer/virksomheder 
    \item Ikke-omgående aflivning → kun hos bemyndiget
  \end{itemize}
  \item Miljøstyrelsen kan tilbagekalde bemyndigelse ved overtrædelser eller ved manglende formål
  \item \textbf{Budskab:} Streng kontrol – almindelige jagttegnsindehaver må ikke pleje disse arter
\end{itemize}
\end{frame}

\begin{frame}{Bilag 6 – Arter der kun må plejes af bemyndiget person/virksomhed}
\tiny
Tilskadekommet vildt eller nødstedt yngel (§ 30, stk. 3) – må IKKE plejes af private uden bemyndigelse fra Miljøstyrelsen/Styrelsen for Grøn Arealomlægning og Vandmiljø

\begin{columns}[T]
\column{0.5\textwidth}
\textbf{Pattedyr}
\begin{itemize}
  \item Bæver (\textit{Castor fiber})
  \item Flagermus (\textit{Chiroptera spp.}), alle arter
  \item Guldsjakal (\textit{Canis aureus})
  \item Hvaler (\textit{Cetacea spp.}), alle arter
  \item Kalve af krondyr, dådyr, sika (\textit{Cervus/Dama spp.})
  \item Lam af rådyr (\textit{Capreolus capreolus})
  \item Odder (\textit{Lutra lutra})
  \item Pindsvin (\textit{Erinaceus europaeus})
  \item Sæler (\textit{Phocidae spp.}), alle arter
  \item Ulv (\textit{Canis lupus})
\end{itemize}

\column{0.5\textwidth}
\textbf{Fugle}
\begin{itemize}
  \item Biæder (\textit{Merops apiaster})
  \item Ellekrage (\textit{Coracias garrulus})
  \item Engsnarre (\textit{Crex crex})
  \item Hærfugl (\textit{Upupa epops})
  \item Høgesanger (\textit{Curruca nisoria})
  \item Hvid stork (\textit{Ciconia ciconia})
  \item Hvidbrystet præstekrave (\textit{Charadrius alexandrinus})
  \item Isfugl (\textit{Alcedo atthis})
  \item Rørdrum (\textit{Botaurus stellaris})
  \item Rovfugle (\textit{Accipitriformes spp.}), alle arter
  \item Rovterne (\textit{Sterna caspia})
  \item Sandterne (\textit{Sterna nilotica})
  \item Sort stork (\textit{Ciconia nigra})
  \item Sortterne (\textit{Chlidonias niger})
  \item Trane (\textit{Grus grus})
  \item Ugler (\textit{Strigiformes spp.}), alle arter
\end{itemize}
\end{columns}

\tiny
Andre arter (ikke på bilag 6): Må plejes af alle med henblik på genudsætning (§ 29). Kontakt Styrelsen for Grøn Arealomlægning og Vandmiljø ved tvivl.
\end{frame}

